% =========================================================
% Volume IV  \S{}1.7.?  --- N as a Model of the Arithmetic Signature
% WIRING NOTE (cross-volume):
%   Verify the exact Volume II labels before final commit. Intended targets:
%   def:arithmetic-signature, def:peano-system,
%   def:addition-on-peano-system, def:multiplication-on-peano-system,
%   thm:n-is-peano-system, thm:n-is-commutative-semiring,
%   thm:n-is-well-ordered, thm:n-embeds-in-z.
% =========================================================

\subsection*{The Natural Numbers as a Structure}

\begin{remark*}[Modeling Goal]
We exhibit $\mathbb{N}$ as a concrete \emph{model}: fix the arithmetic
signature, take the carrier and operations built in Volume~II, interpret each
symbol, and discharge satisfaction by citing the Volume~II theorem that
$\mathbb{N}$ obeys the relevant natural-number axioms. The successor fragment
recovers the Peano-system view; the additive and multiplicative fragments are
obtained by reduction.
\end{remark*}

% ---- 1. the signature being interpreted ----
\begin{remark*}[Exposition]
A model is a signature together with an interpretation on a carrier. The
signature here is the natural-number fragment of the
\hyperref[def:arithmetic-signature]{arithmetic signature},
\[
\mathcal{L}_{\mathbb{N}}=\{S,+,\cdot,0,1\},
\]
with unary function symbol $S$, binary function symbols $+,\cdot$, and
constant symbols $0,1$. The order symbol $<$ is deliberately absent; it enters
by expansion below.
\end{remark*}

% ---- 2. the model itself ----
\begin{definition}[The Natural Number Model $\mathcal{N}$]
\label{def:natural-number-model}
The \textbf{natural number model} is the $\mathcal{L}_{\mathbb{N}}$-structure
\[
\mathcal{N}
=
\bigl(\mathbb{N};\,
S^{\mathcal{N}},\,+^{\mathcal{N}},\,\cdot^{\mathcal{N}},
\,0^{\mathcal{N}},\,1^{\mathcal{N}}\bigr),
\]
whose carrier $|\mathcal{N}|=\mathbb{N}$ is the set of natural numbers
constructed in Volume~II, and whose symbols are interpreted by the
\hyperref[def:addition-on-peano-system]{natural-number addition} and
\hyperref[def:multiplication-on-peano-system]{natural-number multiplication}:
$S^{\mathcal{N}}$ is successor, $+^{\mathcal{N}}$ is natural-number addition,
$\cdot^{\mathcal{N}}$ is natural-number multiplication, and
$0^{\mathcal{N}},1^{\mathcal{N}}$ are the natural-number constants $0$ and $1$.
\end{definition}

\begin{remark*}[Standard quantified statement]
\[
|\mathcal{N}|=\mathbb{N},\qquad
S^{\mathcal{N}}:\mathbb{N}\to\mathbb{N},\qquad
+^{\mathcal{N}},\cdot^{\mathcal{N}}:\mathbb{N}\times\mathbb{N}\to\mathbb{N},
\qquad
0^{\mathcal{N}},1^{\mathcal{N}}\in\mathbb{N}.
\]
\end{remark*}\begin{remark*}[Interpretation]
Nothing here is new mathematics: the carrier, successor, addition, and
multiplication are exactly the objects built in Volume~II. The model only
packages them as a single object $\mathcal{N}$, so that ``$\mathbb{N}$'' can be
handed to any theorem that quantifies over arithmetic structures. The signature
is the type; $\mathcal{N}$ is the instance.
\end{remark*}

% ---- 3. satisfaction: the keystone, wired to Volume II ----
\begin{remark*}[Satisfaction]
$\mathcal{N}$ is a \emph{model} of the natural-number axioms:
\[
\mathcal{N}\models \Sigma_{\mathbb{N}},
\qquad\text{equivalently}\qquad
\operatorname{PeanoSystem}(\mathbb{N},S,0).
\]
This is not re-proved here. It is exactly the content of
\hyperref[thm:n-is-peano-system]{$\mathbb{N}$ Is a Peano System} in
Volume~II, which verifies the Peano axioms against the constructed successor
operation. The algebraic arithmetic fragment is witnessed separately by
\hyperref[thm:n-is-commutative-semiring]{$\mathbb{N}$ Is a Commutative Semiring},
which verifies the laws of addition and multiplication on $\mathbb{N}$.
\end{remark*}

% ---- 4. expansion to the ordered natural-number model ----
\begin{remark*}[Expansion to the ordered natural-number model]
Adjoining the relation symbol $<$ gives the expansion
\[
\mathcal{N}_{<}=(\mathbb{N};S,+,\cdot,0,1,<),
\qquad
\operatorname{Expansion}(\mathcal{N}_{<},\mathcal{N}),
\]
in the sense of \hyperref[def:structure-expansion]{Structure Expansion}: the
carrier is unchanged and exactly one new interpreted symbol is added. Here
$<^{\mathcal{N}}$ is the usual order on $\mathbb{N}$. That this expansion has
the expected order behavior is witnessed by
\hyperref[thm:n-is-well-ordered]{$\mathbb{N}$ Is Well-Ordered} in Volume~II.
\end{remark*}

% ---- 5. reduct to the Peano successor structure ----
\begin{remark*}[Reduct to the successor structure]
Forgetting $+,\cdot,1$ leaves the successor reduct
\[
\mathcal{N}\!\restriction\!\{S,0\}
=
(\mathbb{N};S,0),
\qquad
\operatorname{Reduction}\!\bigl(\mathcal{N}\!\restriction\!\{S,0\},\,\mathcal{N}\bigr),
\]
in the sense of \hyperref[def:structure-reduction]{Structure Reduction}: the
same carrier with fewer visible symbols. This reduct is the Peano-system view
of the natural numbers.
\end{remark*}

% ---- 6. reduct to the additive monoid ----
\begin{remark*}[Reduct to the additive monoid]
Forgetting $S,\cdot,1$ leaves the additive reduct
\[
\mathcal{N}\!\restriction\!\{+,0\}
=
(\mathbb{N};+,0),
\qquad
\operatorname{Reduction}\!\bigl(\mathcal{N}\!\restriction\!\{+,0\},\,\mathcal{N}\bigr).
\]
This reduct is the additive commutative monoid of natural numbers: it remembers
addition and the additive identity, but forgets multiplication and successor as
primitive symbols.
\end{remark*}

% ---- 7. reduct to the multiplicative monoid ----
\begin{remark*}[Reduct to the multiplicative monoid]
Forgetting $S,+,0$ leaves the multiplicative reduct
\[
\mathcal{N}\!\restriction\!\{\cdot,1\}
=
(\mathbb{N};\cdot,1),
\qquad
\operatorname{Reduction}\!\bigl(\mathcal{N}\!\restriction\!\{\cdot,1\},\,\mathcal{N}\bigr).
\]
This reduct is the multiplicative commutative monoid of natural numbers.
\end{remark*}

% ---- 8. blueprint position + embedding chain ----
\begin{remark*}[Blueprint position]
In the structural blueprint $\mathcal{N}$ is the initial counting structure: it is
the first model in the number-system hierarchy and the source of the successor,
addition, multiplication, and induction principles used later. It sits at the head
of the canonical embedding chain
\[
\mathbb{N}\hookrightarrow\mathbb{Z}\hookrightarrow\mathbb{Q}\hookrightarrow\mathbb{R}.
\]
The first arrow is the canonical embedding of Volume~II,
\hyperref[thm:n-embeds-in-z]{$\mathbb{N}$ Embeds in $\mathbb{Z}$}. Passing
from $\mathbb{N}$ to $\mathbb{Z}$ supplies additive inverses; passing from
$\mathbb{Z}$ to $\mathbb{Q}$ supplies multiplicative inverses for nonzero
elements; passing from $\mathbb{Q}$ to $\mathbb{R}$ supplies the analytic
completeness needed for real analysis.
\end{remark*}

\begin{remark*}[Interpretation]
$\mathbb{N}$ is not a group, ring, or field under its usual arithmetic operations:
subtraction and division are not closed operations on all of $\mathbb{N}$. Its
natural algebraic home is the Peano-system and commutative-semiring layer.
This is exactly why the number-system hierarchy proceeds from
$\mathbb{N}$ to $\mathbb{Z}$, then to $\mathbb{Q}$, and finally to
$\mathbb{R}$.
\end{remark*}

\begin{dependencies}
\begin{itemize}
  \item \hyperref[def:arithmetic-signature]{Arithmetic Signature}.
  \item \hyperref[def:peano-system]{Peano System}.
  \item \hyperref[def:natural-numbers]{The Natural Numbers}.
  \item \hyperref[def:addition-on-peano-system]{Addition on $\mathbb N$}.
  \item \hyperref[def:multiplication-on-peano-system]{Multiplication on $\mathbb N$}.
\end{itemize}
\end{dependencies}
